% Section 2.3, from lectures/L02/README.md: the objectives, the evaluation questions, and the next
% lecture.

\needspace{10\baselineskip}
\section{Review}
\label{c2:sec:review}

\subsection*{What you should be able to do now}
After this chapter, you should:
\begin{itemize}
  \item Be able to randomize the order of the training sets before each new epoch.
  \item Be able to compute the precision of a trained regression model.
  \item Be able to determine when the model is sufficiently trained based on a given threshold
    value.
  \item Be able to adjust the learning rate during training, based on how much the model is
    improving.
\end{itemize}

\needspace{6\baselineskip}
\subsection*{Questions to test yourself}
You should be able to explain:
\begin{enumerate}
  \item Why is it beneficial to randomize the training order every epoch?
  \item Why is a static local variable used in \code{initRandGen()} to ensure the random number
    generator is only initialized once? Why does the function belong in \file{ml/utils.cpp} rather
    than in an anonymous namespace in each model's source file?
  \item How is the precision of a linear regression model computed?
  \item What does it mean for the precision to exceed a given threshold, and how is this used to
    stop training early?
  \item Why is precision calculation important when training machine learning models?
  \item How does \code{ml::lin_reg::Adaptive} decide whether to raise or lower its learning rate,
    and why is the rate clamped at both ends?
\end{enumerate}

\needspace{6\baselineskip}
\subsection*{Looking ahead}
Chapter~\ref{c3:ch:nn} is about neural networks:
\begin{itemize}
  \item An introduction to neural networks: feedforward, backpropagation, gradient descent, and
    activation functions.
  \item Training neural networks by hand.
  \item The implementation of a dense layer interface and a stub.
\end{itemize}
